\begin{table*}[htbp]
\centering
\small
\caption{Validation strategies and evidence strength.}
\label{tab:validation}
\begin{tabularx}{\textwidth}{@{}>{\raggedright\arraybackslash}p{0.20\textwidth}YY@{}}
\toprule
\textbf{Validation dimension} & \textbf{Current extracted signal} & \textbf{Interpretation} \\
\midrule
Validation approaches & Empirical / patient-data validation (25); Simulation / benchmark comparison (18); Sensitivity / uncertainty analysis (13); Cross-validation / holdout (8); Conceptual / no model evaluation (6); Expert evaluation (3) & Reports record-level mentions of empirical, simulated, cross-validation, uncertainty, expert, and conceptual approaches. \\
Validation data source & Patient/clinical, simulation/benchmark, and other sources remain free-text and are preserved in the record-level output. & Avoids imposing an unsupported pooled effect interpretation. \\
Fidelity and uncertainty & Sensitivity/uncertainty language appears in 13 records. & Shows how often model credibility was discussed explicitly. \\
Reproducibility & Explicit affirmative, code, open-source, or reproducibility language appears in 26 records. & Identifies an actionable transparency gap. \\
Clinical readiness & Prototype / proof of concept (43); Theoretical / conceptual (21); Deployed / operational (19); Prospective / clinical study (3); Retrospective validation (1) & Separates conceptual promise, prototypes, validation studies, and deployed systems. \\
\bottomrule
\end{tabularx}
\end{table*}
